%%%%%%%%%%  scan idx 44  |  folio 31  %%%%%%%%%%
\SGAPrintPage{31}%
Formula
(1.22) admits the following generalizations: let $K$ be any $\lambda$-ring, otherwise arbitrary; for $x\in K$, put
\sgaStableTarget{sga6.e00.eq.1-23}\begin{equation}\tag{1.23}
\gamma^n(x)=(-1)^n\sum_{i=0}^n(-1)^i\lambda^i(x+n)=\lambda^n(x+n-1)
\end{equation}
(the identity of the last two members follows from an immediate induction). In particular, we have $\gamma^0(x)=1$ and $\gamma^1(x)=x$; next put:
\sgaStableTarget{sga6.e00.eq.1-24}\begin{equation}\tag{1.24}
\gamma_t(x)=\sum_{n\geq0}\gamma^n(x)t^n\in K[[t]]
\end{equation}
whence we deduce the formula:
\sgaStableTarget{sga6.e00.eq.1-25}\begin{equation}\tag{1.25}
\gamma_t(x)=\lambda_{t/(1-t)}(x),
\end{equation}
which is equivalent to
\sgaStableTarget{sga6.e00.eq.1-25bis}\begin{equation}\tag{1.25 bis}
\lambda_s(x)=\gamma_{s/(1+s)}(x).
\end{equation}
This proves that the $\lambda^i$ are expressed by means of the $\gamma^i$ and that $\gamma_t(x+y)=\gamma_t(x)\gamma_t(y)$, and consequently that the maps $\gamma^i$ define a new $\lambda$-ring structure on $K$. Formula (1.22) then shows that, for $K=\widetilde A$ and every $x\in\widetilde A$, we have
\sgaStableTarget{sga6.e00.eq.1-22bis}\begin{equation}\tag{1.22 bis}
\gamma^n(x-\varepsilon(x))=[0,1+(-1)^{n-1}(n-1)!x^n+\cdots]
\end{equation}
and consequently, the \emph{component} $x^n$ \emph{is obtained, up to the factor} $(-1)^{n-1}(n-1)!$, \emph{by reducing modulo} $\widetilde A^{n+1}$ \emph{the element} $\gamma^n(x-\varepsilon(x))$ \emph{of} $\widetilde A^n$.

\sgaStableTarget{sga6.e00.the-chern-homomorphism}\paragraph{The Chern homomorphism:}
For every graded ring $A$ satisfying the preceding conditions, we are going to define a homomorphism of augmented rings:
\sgaStableTarget{sga6.e00.eq.1-26}\begin{equation}\tag{1.26}
\ch:\widetilde A\longrightarrow \widehat{A\otimes\Q}.
\end{equation}
This homomorphism is defined unambiguously by the conditions that it be an additive homomorphism, take $1$ to $1$, be such that the com%
%%%%%%%%%%  scan idx 45  |  folio 32  %%%%%%%%%%
\SGAPrintPage{32}%
ponents of $\ch(x)$ in degree $>0$ be given by universal isobaric polynomials with rational coefficients in the components of $x$, and that
\sgaStableTarget{sga6.e00.eq.1-27}\begin{equation}\tag{1.27}
\ch([1,1+x^1])=\exp x^1=\sum_{n\geq0}(x^1)^n/n!.
\end{equation}
We then immediately obtain the explicit expression for $\ch$: let $\eta$ be the additive endomorphism of $\widehat{A\otimes\Q}$ defined, for $x^i$ of degree $i$, by
\sgaStableTarget{sga6.e00.eq.1-28}\begin{equation}\tag{1.28}
\eta(x^i)=1/(-1)^{i-1}(i-1)!\,x^i
\end{equation}
we then have
\sgaStableTarget{sga6.e00.eq.1-29}\begin{equation}\tag{1.29}
\ch\left([n,1+\sum_{i\geq1}x^i]\right)=n+\eta\left(\log\left(1+\sum_{i\geq1}x^i\right)\right).
\end{equation}
As we can ``invert'' this formula, at least when $A$ has only finitely many homogeneous components, we see that, in this case, we have an \emph{isomorphism} $\ch$ \emph{from the ring} $\widetilde A\otimes\Q$ \emph{onto the ring} $\widehat{A\otimes\Q}$ (which completely elucidates the structure of $\widetilde A\otimes\Q$).

Recall that, according to Hirzebruch [1], a formal series $f\in\Q[[t]]$ defines, by universal polynomial formulas, an additive homomorphism $1+\Ahp\longrightarrow 1+(\widehat{A\otimes\Q})^+$, denoted $\mathcal T_f$; when $f(t)=t/(1-\exp(-t))$, we simply put $\mathcal T_f=\mathcal T$ ($\mathcal T$ is the initial of Todd). We extend the definition of $\mathcal T(x)$ to $x\in\widetilde A$. This said, we have the following result:

\sgaStableTarget{sga6.e00.prop.1-3}\textbf{Proposition 1.3.} \emph{Let $N\in\widetilde A$ be of the form $[q,1+N^1+\cdots+N^q]$} (we have $N^i=0$ for $i>q$). \emph{If we put:}
\[
\lambda_{-1}(N)=\sum_{i\geq0}(-1)^i\lambda^i(N)=\sum_{i=0}^{q}(-1)^i\lambda^i(N),
\]
\emph{we shall have:}
\sgaStableTarget{sga6.e00.eq.1-30}\begin{equation}\tag{1.30}
\ch(\lambda_{-1}(N))=(-1)^q N^q\mathcal T(\check N)^{-1},
\end{equation}
%%%%%%%%%%  scan idx 46  |  folio 33  %%%%%%%%%%
\SGAPrintPage{33}%
\emph{which is equivalent to:}
\sgaStableTarget{sga6.e00.eq.1-30bis}\begin{equation}\tag{1.30 bis}
\ch(\lambda_{-1}(\check N))=N^q\mathcal T(N)^{-1}.
\end{equation}

\sgaStableTarget{sga6.e00.sec.4.occ2}\subsection*{\S 4. The operations $\lambda^p(N,x)$}

Let $A$ be the ring of formal series, with coefficients in $\Z$, in two infinite sequences of indeterminates, denoted $\lambda^iN$ and $\lambda^i x$ ($i\geq1$). Put $\lambda^1N=N$ and $\lambda^1x=x$, and introduce on $A$ the unique special $\lambda$-ring structure for which we have $\lambda^i(N)=\lambda^iN$ and $\lambda^i(x)=\lambda^ix$ and for which the operations $\lambda^i$ are ``continuous'' (if we restricted ourselves to \emph{polynomials} in indeterminates $\lambda^iN$ and $\lambda^ix$, we would have the ``free $\lambda$-ring generated by $N$ and $x$''; our ring $A$ is a completion of it). We can then form
\sgaStableTarget{sga6.e00.eq.1-31}\begin{equation}\tag{1.31}
\lambda_{-1}(N)=\sum_{i\geq0}(-1)^i\lambda^iN,
\end{equation}
which is a formal series with constant term $1$, hence invertible; we can therefore define an element $\lambda^p(N,x)$ of $A$ unambiguously by putting
\sgaStableTarget{sga6.e00.eq.1-32}\begin{equation}\tag{1.32}
\lambda^p(N,x)\lambda_{-1}(N)=\lambda^p(x\lambda_{-1}(N))
\end{equation}

\sgaStableTarget{sga6.e00.thm.1-4}\textbf{Theorem 1.4.} \emph{Let $J_q$ be the closed ideal of $A$ generated by the $\lambda^iN$ with $i>q$; identify $A/J_q$ with the ring of formal series in the $\lambda^i x$ and the $\lambda^jN$ ($1\leq j\leq q$). Then $\lambda^p(N,x)$ reduced modulo $J_q$ is a} polynomial, \emph{isobaric of weight $p$ with respect to the variables $\lambda^i x$.}

\emph{Let $N$ and $F$ be two vector spaces of finite dimensions $q,q'$ over an algebraically closed field $k$ of characteristic $0$; let $N^{(p)}$ be the kernel of the map} $(a_1,\ldots,a_p)\longrightarrow a_1+\cdots+a_p$ from $N^p$ \emph{to} $N$.
We
%%%%%%%%%%  scan idx 47  |  folio 34  %%%%%%%%%%
\SGAPrintPage{34}%
\emph{let the symmetric group $\mathfrak S_p$ operate on $N^p$, hence on $N^{(p)}$ hence on $\Lambda(N^{(p)})$, then on $\bigotimes^p F$ and on $\Lambda(N^{(p)})\otimes\bigotimes^p F$. Consider this vector space as a graded module over $\mathfrak S_p\times\Gl(N)\times\Gl(F)$ and take its ``alternating part'', which is a graded module over $\Gl(N)\times\Gl(F)\eqqcolon G$. We then have}
\sgaStableTarget{sga6.e00.eq.1-33}\begin{equation}\tag{1.33}
\lambda^p(N,F)=\underline{E}_G\left(\left(\Lambda(N^{(p)})\otimes\bigotimes^p F\right)^{\mathrm{alt}}\right).
\end{equation}

N.B. If $M$ is a complex of $G$-modules, we denote by $\underline{E}_G(M)$ the alternating sum in $\underline{K}(G)$ of the classes of the components of $M$; this is the Euler--Poincaré characteristic of $M$; in formula (1.33), the first member $\lambda^p(N,F)$ denotes the element of the ring $\underline{K}(G)$ obtained by substituting for the variables $\lambda^iN$ and $\lambda^iF$ the classes in $\underline{K}(G)$ of the $G$-modules $\Lambda^i(N)$ and $\Lambda^i(F)$ respectively.

\textbf{Proof:} By expanding the second member of (1.32) by the second formula (1.11), we see that the first assertion is proved if we prove the analogous assertion obtained by putting $x=1$, i.e. if we prove that, in the $\lambda$-ring $\Z[\lambda^1N,\ldots,\lambda^qN]$ where $\lambda^iN=0$ for $i>q$, the element $\lambda^p(\lambda_{-1}(N))$ is divisible by $\lambda_{-1}(N)$. A priori, the quotient $\lambda^p(N,1)=\lambda^p(\lambda_{-1}(N))/\lambda_{-1}(N)$ is a formal series contained in the field of fractions $K$ of the ring of formal series in the $\lambda^iN$ for $1\leq i\leq q$. On the other hand, also denote by $N$ a vector space of dimension $q$ over an algebraically closed field $k$ of characteristic $0$, and consider the element $\ell^p(N)=\underline{E}_G((\Lambda(N^{(p)}))^{\mathrm{alt}})\in \underline{K}(G)$, where we now have $G=\Gl(N)$. We know (\S 2, theorem 1.1) that $\underline{K}(G)$ is identified
%%%%%%%%%%  scan idx 48  |  folio 35  %%%%%%%%%%
\SGAPrintPage{35}%
with the
subring of $K$ generated by the $\lambda^iN$ for $1\leq i\leq q$ and by the inverse of $\lambda^qN$. On the other hand, we readily verify in this ring that we have $\ell^p(N)\lambda_{-1}(N)=\lambda^p(\lambda_{-1}(N))$. As we are in a large field $K$ and $\lambda_{-1}(N)\neq0$, we conclude that $\ell^p(N)=\lambda^p(N,1)$; both members are in the intersection of $\underline{K}(G)$ and the ring of formal series in the $\lambda^iN$, and are therefore polynomials in the $\lambda^iN$.

To prove (1.33) in the general case, we verify that the second member satisfies in $\underline{K}(G)$ the formula analogous to (1.32), that is, its product by $\lambda_{-1}(N)$ is $\lambda^p(F\lambda_{-1}(N))$. Likewise, the product of $\lambda^p(N,F)\in \underline{K}(G)$ by $\lambda_{-1}(N)$ is $\lambda^p(F\lambda_{-1}(N))$, as follows from the fact that $\underline{K}(G)$ is \emph{special}. As $\underline{K}(G)$ is an integral domain (cf.\ \S 2) and $\lambda_{-1}(N)\neq0$ (same reference), we have indeed proved formula (1.33). q.e.d.

\sgaStableTarget{sga6.e00.rem.unnamed.occ2}\textbf{Remarks:} a) When the dimension $q'$ of $F$ is $\geq p$, then (1.33) characterizes the polynomial $\lambda^p(N,x)$ in the $\lambda^iN$ ($1\leq i\leq q$) and the $\lambda^jx$ ($1\leq j\leq p$).

b) It is plausible that (1.33) still holds when the field $k$ is not algebraically closed\SGAAuthorFootnote{*}{This is indeed proved by the same method.}. This question would have to be resolved to prove the Riemann--Roch theorem over a field of characteristic $0$ that is not algebraically closed by the method we are going to explain. Moreover, the passage from an $\mathfrak S_p$-module to its alternating part has a reasonable meaning only when the base field has characteristic $0$.

Now let $K$ be any \emph{special} $\lambda$-ring and let $N$ be an
ele%
